\begin{python0}
from solutions import *; clear()
\end{python0}
\sectionthree{Private members of the parent class}

For public inheritance, a child cannot access private member variables
or private methods of parent. Let's do an experiment to verify this.
Here's our Basic Example:
\begin{console}
#include <iostream>

class P
{
public:
        P() : x_(1) {}
        int get_x() const { return x_; }
        void set_x(int x) { x_ = x; }
private:
        int x_;
};

class C: public P
{
public:
        C() : x_(3), y_(2) {}
        int get_y() const { return y_; }
        void set_y(int y) { y_ = y; }
private:
        int x_;
        int y_;
};

int main()
{   
    C c;
    std::cout << "c.x_ =" << c.get_x() << ", "
              << "c.y_ =" << c.get_y() << '\n';
    std::cout << "c.x_ =" << c.P::get_x() << ", "
              << "c.y_ =" << c.C::get_y() << '\n';
return 0;
}
\end{console}

Now let's create a method in \verb!C! to access a private member of
\verb!P!:
\begin{console}[commandchars=\~\@\$]
...
class C : public P
{
public:
        ...
        ~EMPHASIZE@~redtext@void m() { P::x_ = 42; }$$
private:
        int x_;
        int y_;
};

... 
\end{console}

You will get an error when you compile your program -- make sure you
read the error message. This is the same for private method in the
parent:
\begin{console}[commandchars=\~\@\$]
...

class P
{     
...
private:
        ~EMPHASIZE@~redtext@void m() {}$$
        int x_;
};

class C: public P
{
public:
        ...
        ~EMPHASIZE@~redtext@void n() { m(); }$$
        ...
};

...
\end{console}

